% GENERATED FILE. Edit canonical lessons, then run npm run generate:books.
% canonical-source-hash: fnv1a64:078be273a6adc6a2
% canonical-lessons: PT-C352-reacao, PT-C352-reagir, PT-C352-pretexto, PT-C352-desculpa, PT-C352-justificar

\chapter{Como reagir}
\label{ch:pt-como-reagir}
\hlchaptermodality{352}

\begin{chapteropening}
\textbf{By the end of this chapter:} \emph{I can describe how people react, give an excuse or a pretext, and justify a decision.}

Complete the last lesson of chapter 352: I can describe how people react, give an excuse or a pretext, and justify a decision.
\end{chapteropening}

\section[a reação]{a reação — the reaction}

\label{lesson:PT-C352-reacao}



\begin{quote}
Before the new one: say the Portuguese for to entail, then the Portuguese for to trigger.
\end{quote}

\subsection*{You'll want to know: a reação}
\textbf{a reação} — \textquotedblleft{}the reaction\textquotedblright{}.

\begin{quote}
A reação dele foi uma surpresa para todos.
\end{quote}

\begin{quote}
\emph{His reaction was a surprise to everyone.}
\end{quote}

AO90 \emph{reação}; older \emph{reacção}.

\begin{cousinweb}
\emph{reação} is a learned word modelled on French \emph{réaction} (EtymDB), from Medieval Latin \emph{reactiō}, built on \emph{re-} \textquotedblleft{}back\textquotedblright{} and \emph{agere} \textquotedblleft{}to do, to act\textquotedblright{} (Lewis \& Short).
\end{cousinweb}

\begin{grammarlens}[title={What you've built}]
One more word for telling, explaining and describing at B1.
\end{grammarlens}

\subsection*{Your turn}
\begin{itemize}
\raggedright
  \item \emph{Say it:} \emph{a reação}
  \item \emph{Say it:} \emph{a reação}, once more
  \item \emph{Say it:} a sentence of your own with \emph{a reação}
  \item \emph{Recall:} say the Portuguese for to entail, then the Portuguese for to trigger, then say \emph{a reação} again
\end{itemize}

\begin{tcolorbox}[breakable,colback=teal!4,colframe=teal!35!black,title={Before you move on}]
What does a reação mean? (\textquotedblleft{}The reaction\textquotedblright{}.) Say it once more.
\end{tcolorbox}
\section[reagir]{reagir — to react}

\label{lesson:PT-C352-reagir}



\begin{quote}
Before the new one: say the Portuguese for to trigger, then the Portuguese for the reaction.
\end{quote}

\subsection*{You'll want to know: reagir}
\textbf{reagir} — \textquotedblleft{}to react\textquotedblright{}.

\begin{quote}
Ele reagiu mal à notícia.
\end{quote}

\begin{quote}
\emph{He reacted badly to the news.}
\end{quote}

\emph{Reagir a} — to react to. Spelling: \emph{reajo, reages, reage}.

\begin{cousinweb}
\emph{reagir} is a learned formation on Latin \emph{re-} \textquotedblleft{}back\textquotedblright{} and \emph{agere} \textquotedblleft{}to do, to act\textquotedblright{} (EtymDB, Lewis \& Short).
\end{cousinweb}

\begin{grammarlens}[title={What you've built}]
One more word for telling, explaining and describing at B1.
\end{grammarlens}

\subsection*{Your turn}
\begin{itemize}
\raggedright
  \item \emph{Say it:} \emph{reagir}
  \item \emph{Say it:} \emph{reagir}, once more
  \item \emph{Say it:} a sentence of your own with \emph{reagir}
  \item \emph{Recall:} say the Portuguese for to trigger, then the Portuguese for the reaction, then say \emph{reagir} again
\end{itemize}

\begin{tcolorbox}[breakable,colback=teal!4,colframe=teal!35!black,title={Before you move on}]
What does reagir mean? (\textquotedblleft{}To react\textquotedblright{}.) Say it once more.
\end{tcolorbox}
\section[o pretexto]{o pretexto — the pretext}

\label{lesson:PT-C352-pretexto}



\begin{quote}
Before the new one: say the Portuguese for the reaction, then the Portuguese for to react.
\end{quote}

\subsection*{You'll want to know: o pretexto}
\textbf{o pretexto} — \textquotedblleft{}the pretext\textquotedblright{}.

\begin{quote}
Ela usou o trabalho como pretexto para não ir.
\end{quote}

\begin{quote}
\emph{She used work as a pretext for not going.}
\end{quote}

\emph{Sob o pretexto de} — on the pretext of.

\begin{cousinweb}
\emph{pretexto} is a learned borrowing from Latin \emph{praetextus} \textquotedblleft{}an outward appearance, a show\textquotedblright{} (EtymDB, Lewis \& Short), from \emph{praetexere} \textquotedblleft{}to weave in front\textquotedblright{}, built on \emph{texere} \textquotedblleft{}to weave\textquotedblright{}.
\end{cousinweb}

\begin{grammarlens}[title={What you've built}]
One more word for telling, explaining and describing at B1.
\end{grammarlens}

\subsection*{Your turn}
\begin{itemize}
\raggedright
  \item \emph{Say it:} \emph{o pretexto}
  \item \emph{Say it:} \emph{o pretexto}, once more
  \item \emph{Say it:} a sentence of your own with \emph{o pretexto}
  \item \emph{Recall:} say the Portuguese for the reaction, then the Portuguese for to react, then say \emph{o pretexto} again
\end{itemize}

\begin{tcolorbox}[breakable,colback=teal!4,colframe=teal!35!black,title={Before you move on}]
What does o pretexto mean? (\textquotedblleft{}The pretext\textquotedblright{}.) Say it once more.
\end{tcolorbox}
\section[a desculpa]{a desculpa — the excuse; the apology}

\label{lesson:PT-C352-desculpa}



\begin{quote}
Before the new one: say the Portuguese for to react, then the Portuguese for the pretext.
\end{quote}

\subsection*{You'll want to know: a desculpa}
\textbf{a desculpa} — \textquotedblleft{}the excuse; the apology\textquotedblright{}.

\begin{quote}
Ele chegou tarde e deu uma desculpa estranha.
\end{quote}

\begin{quote}
\emph{He arrived late and gave a strange excuse.}
\end{quote}

\emph{Pedir desculpa} — to apologise; \emph{desculpa!} — sorry!

\begin{cousinweb}
\emph{desculpa} is built on the verb \emph{desculpar} \textquotedblleft{}to excuse\textquotedblright{}, from \emph{des-} and \emph{culpa} \textquotedblleft{}fault, blame\textquotedblright{}, from Latin \emph{culpa} \textquotedblleft{}fault\textquotedblright{} (Lewis \& Short).
\end{cousinweb}

\begin{grammarlens}[title={What you've built}]
One more word for telling, explaining and describing at B1.
\end{grammarlens}

\subsection*{Your turn}
\begin{itemize}
\raggedright
  \item \emph{Say it:} \emph{a desculpa}
  \item \emph{Say it:} \emph{a desculpa}, once more
  \item \emph{Say it:} a sentence of your own with \emph{a desculpa}
  \item \emph{Recall:} say the Portuguese for to react, then the Portuguese for the pretext, then say \emph{a desculpa} again
\end{itemize}

\begin{tcolorbox}[breakable,colback=teal!4,colframe=teal!35!black,title={Before you move on}]
What does a desculpa mean? (\textquotedblleft{}The excuse; the apology\textquotedblright{}.) Say it once more.
\end{tcolorbox}
\section[justificar]{justificar — to justify}

\label{lesson:PT-C352-justificar}



\begin{quote}
Before the new one: say the Portuguese for the pretext, then the Portuguese for the excuse.
\end{quote}

\subsection*{You'll want to know: justificar}
\textbf{justificar} — \textquotedblleft{}to justify\textquotedblright{}.

\begin{quote}
Ele tentou justificar a decisão.
\end{quote}

\begin{quote}
\emph{He tried to justify the decision.}
\end{quote}

\emph{Justificar-se} — to explain yourself.

\begin{cousinweb}
\emph{justificar} is a learned borrowing from Late Latin \emph{iūstificāre} \textquotedblleft{}to act justly, to justify\textquotedblright{} (EtymDB), built on \emph{iūstus} \textquotedblleft{}fair, righteous\textquotedblright{} and \emph{facere} \textquotedblleft{}to make\textquotedblright{} (Lewis \& Short).
\end{cousinweb}

\begin{grammarlens}[title={What you've built}]
One more word for telling, explaining and describing at B1.
\end{grammarlens}

\subsection*{Your turn}
\begin{itemize}
\raggedright
  \item \emph{Say it:} \emph{justificar}
  \item \emph{Say it:} \emph{justificar}, once more
  \item \emph{Say it:} a sentence of your own with \emph{justificar}
  \item \emph{Recall:} say the Portuguese for the pretext, then the Portuguese for the excuse, then say \emph{justificar} again
\end{itemize}

\begin{tcolorbox}[breakable,colback=teal!4,colframe=teal!35!black,title={Before you move on}]
What does justificar mean? (\textquotedblleft{}To justify\textquotedblright{}.) Say it once more.
\end{tcolorbox}
